% =====================================================================
% NeurIPS 2026 workshop submission (double-blind, non-archival).
% Workshop: "Who Verifies the Agents? Toward Reliable Agent Development"
% Page limit: 4-9 pages of body, EXCLUDING references and appendices.
% Deadline: 2026-08-29, 23:59 AoE. Submission via OpenReview.
%
% Pre-results draft of the self-report study. The INSTRUMENT results
% (Section 3) are measured and real; the ARM results are not yet run.
% Design source:
%     docs/research/neurips-2026-workshop/25-placebo-selfreport-design.md
% Decisions: DL-49..DL-52, DL-77..DL-90 in 15-decision-log.md.
%
% NUMBERS POLICY.
%   * Quantities describing the FROZEN POOL, MODEL SELECTION and the
%     INSTRUMENT (polarity bias, phrasing floor, dissociation) are real and
%     come from build/research/placebo/{failure_pool,preflight}.json and
%     build/research/placebo/instrument/. They may be checked today.
%   * Every quantity describing an ARM OUTCOME does not exist yet. No arm has
%     been run. Those slots are \result{...} and render visibly in red.
%   * When the confirmatory batch lands, each \result{} is replaced by a value
%     emitted by the analysis script, never hand-entered.
%
% PLACEHOLDERS are \TODO{...} (missing prose/asset) and \result{...} (missing
% number). Both must be zero before submission; \draftnotice must be deleted.
%
% PAGE BUDGET. Measured 2026-07-27 against the SUBMISSION form (\draftnotice
% commented out), by locating section starts in the built PDF.
%
%   state            body ends   References   body chars on p10
%   start of day     p10          p11          --
%   after cuts       p10          p10          1584 of 3675  (43%)
%
% One page was recovered. The paper is still OVER by roughly 0.4 page: page 10
% now holds the Conclusion and nothing else, so page 9 is full and the Conclusion
% has nowhere to go.
%
% Cuts taken: the "How this section reads under each outcome" preregistration
% section; a stale Brier-decomposition paragraph in the introduction (it promised
% a primary outcome the paper does not deliver, calibration being unresolvable at
% this n); the self-defect narrative consolidated from four tellings to one plus a
% release note; Related work's three intervention-study paragraphs merged into one
% (priority concession retained, and also stated in the intro); the third
% restatement of "we claim no new decision rule"; the Reproducibility statement
% moved to an appendix, which the venue excludes; and Table 1 replaced by prose --
% it displayed a RETRACTED effect in the paper's most prominent format while the
% following paragraphs establish the spread is not significant (p = 0.178) and the
% ranking does not reproduce (Spearman +0.029).
%
% MEASURED, not assumed: removing \draftnotice does NOT by itself flip the page.
%
% THE REMAINING ~0.4 PAGE IS A JUDGEMENT CALL, NOT A MECHANICAL ONE. Everything
% cheap is spent. The two candidates large enough to close it are:
%     1. The introduction's five-standard enumeration (~1100 chars). This is the
%        paper's thesis setup -- the argument that four mature fields already
%        require this check -- and cutting it weakens the opening.
%     2. The Conclusion's opening framing paragraph (~350 chars) plus further
%        compression of the remaining three. This is the rhetorical payoff.
% Both trade content a reader values for a page. Decide deliberately rather than
% by salami-slicing, which is how the previous 0.6 page was found and is now
% exhausted.
%
% Do NOT trim by dropping the parity, interpretation-licence, or limitations
% material; those are the sections that make the design defensible.
%
% BUILD:  JOB=placebo ./build.sh     or     make JOB=placebo
% =====================================================================

\documentclass{article}

% Style file: neurips_2026.sty, unmodified upstream copy.
% Provenance and SHA-256 are recorded in README.md. Do not edit the .sty.
\usepackage[dblblindworkshop]{neurips_2026}

\workshoptitle{Who Verifies the Agents? Toward Reliable Agent Development}

\usepackage[utf8]{inputenc}
\usepackage[T1]{fontenc}
\usepackage{hyperref}
\usepackage{url}
\usepackage{booktabs}
\usepackage{enumitem}
\usepackage{amsfonts}
\usepackage{amsmath}
\usepackage{nicefrac}
\usepackage{microtype}
\usepackage{xcolor}
\usepackage[normalem]{ulem}
\usepackage{graphicx}     % \sout, for the forbidden-sentence box only

% Placeholder markers. Both render visibly on purpose so they cannot ship
% unnoticed. \TODO is a missing asset or a missing sentence; \result is a
% number that the experiment has not yet produced.
\newcommand{\TODO}[1]{\textcolor{red}{\textbf{[TODO: #1]}}}
\newcommand{\result}[1]{\textcolor{red}{\textbf{[#1]}}}

% Draft banner. DELETE before any build that leaves this machine.
\newcommand{\draftnotice}{%
    \begin{center}
    \fcolorbox{red}{white}{\parbox{0.92\linewidth}{\small\textcolor{red}{%
    \textbf{PRE-RESULTS DRAFT.} No arm of this experiment has been run. The
    frozen failure pool and the model-selection preflight are real and
    reproducible from committed artifacts; every outcome quantity is a
    placeholder of the form \texttt{[NAME]}. This banner and every placeholder
    must be gone before submission.}}}
    \end{center}
}

% Double-blind hygiene: keep identifying strings out of the PDF metadata.
\hypersetup{
    pdfauthor={},
    pdfsubject={},
    pdfkeywords={},
    pdfcreator={},
    pdfproducer={}
}

\title{No Method Passes:\\
A Gauge Study of Agent Self-Report as a Measurement Instrument}

\author{%
    Anonymous Author(s) \\
    Affiliation \\
    Address \\
    \texttt{email} \\
}

\begin{document}

\maketitle

\draftnotice

\begin{abstract}
Verification pipelines for coding agents increasingly read the agent's own
report -- a stated confidence, a request for review -- to decide what a human
must look at. We plant a content effect of known sign in 46 programs -- each verified by
execution, so the mutated variant provably fails a test the original passes
-- hold
everything else fixed, and measure it six times using six semantically identical
wordings of one question. \textbf{The same effect measures $4.17\times$ larger
under one wording than another} ($d = +0.0621$ to $+0.2588$, standard errors
around 0.03; bootstrap 95\% CI on the ratio $[3.1\times, 6.8\times]$, so at least
threefold). Two laboratories running the identical experiment on the identical
programs and reporting the identical estimand would differ several-fold for no
reason but how they asked, and neither would report the phrasing, because no
convention requires it.

Before an instrument is used to detect an effect, four mature fields require it
first be shown capable of resolving one: analytical chemistry has detection
limits, clinimetrics minimal detectable change, clinical trials assay
sensitivity, industrial metrology gauge repeatability and reproducibility, each
with published acceptance thresholds. We run them. Treating wordings as
appraisers and 120 frozen execution-graded programs as parts, we run the same
gauge study over the four elicitation methods in common use and \textbf{all four
are rejected}: \textbf{\%GRR from $37.9\%$ to $61.7\%$} against a rejection
threshold of 30\%, and \textbf{ndc from $1.80$ to $3.45$} resolvable categories
against a required five. The verbalized integer this literature reports most
often is among the weakest, and no format clears either criterion. And
the positive control fails: on a second verified sample, whose floor is measured
over its own parts, a difference planted by construction and confirmed by
execution ($+0.1962$) falls below that instrument's own smallest detectable
change ($0.2776$), and \textbf{no wording resolves it},
so by ICH E10 logic \textbf{a null from this instrument is uninterpretable}
rather than evidence of no effect -- including nulls from the intervention study
our own harness was built to run.

What survives is narrow and we state it as such: the readout \emph{ranks}
correct above incorrect (AUROC 0.717, CI [0.608, 0.819]) while its calibration is
simply unmeasured at this sample size. Even that surviving use is expensive to
leave undocumented -- routing the least-confident $20\%$ of work to a human,
changing nothing but the wording of the question, moves the share of genuine
failures reaching that human from $26.5\%$ to $41.2\%$. And the usual remedy does
not apply: on self-assessed code two unrelated model families are overconfident about the \emph{same}
failing programs ($8$ of $11$ shared, $p = 0.001$) -- three of which pass
\emph{zero} of their $112$, $116$ and $150$ tests -- so a second model is not an
independent second opinion, and agreement is where the danger concentrates. We claim no new decision rule -- the rule
is Currie's, and we name the standard we import. What is new is that the check
has an answer here, and the answer is negative. Elicitation code, frozen
programs, and every artifact are released.
\end{abstract}

% ---------------------------------------------------------------------
\section{Introduction}
% ---------------------------------------------------------------------

A verification pipeline decides what a human must look at. Increasingly it decides
partly by asking the agent: report a confidence, flag work you are unsure about,
escalate when you think you are stuck. Routing, autonomy limits and review budgets
are being built on that answer.

Every field that routes decisions through an instrument first asks whether the
instrument can resolve the thing it is being asked about. Analytical chemistry
will not report a detection below its critical
value~\citep{currie_detection_1995, iso11843_1997}; clinimetrics will not claim a
change smaller than the smallest detectable change~\citep{devet_sdc_2006}; a
clinical trial that cannot recover a known effect is deemed to lack \emph{assay
sensitivity}, and its null is treated as uninterpretable rather than as evidence of
no effect~\citep{ich_e10_2000}; an accredited laboratory must document a decision
rule with a guard band before reporting at all~\citep{iso17025_2017, jcgm106_2012};
and manufacturing rejects a gauge whose measurement variation exceeds thirty per
cent of study variation, or that cannot resolve five distinct
categories~\citep{aiag_msa_2010}. Old, boring, well-specified procedures with
published acceptance thresholds, and they disagree about almost nothing.

None of them has been run on agent self-report. We run them, and the channel
fails: not one of the four elicitation methods in common use passes either
industrial criterion, and the verbalized integer that this literature reports
most often is the weakest of the four. \textbf{We claim no new decision rule --
the rule is Currie's, and we name the standard we import.} What is new is that
the check has an answer here, and that the answer is negative.

Placebo control is the standard fix and, in this setting, no longer novel: three
preregistered, derangement-assigned studies of self-repair in frozen small code
models already exist on exactly these
benchmarks~\cite{iscan_scaffoldnotvocabulary_2026, iscan_falsificationnotexposure_2026,
iscan_formnotcontent_2026}. We claim priority over none of it, and say so here
rather than in a related-work section on page four. What those studies measure is
whether the agent still \emph{succeeds}; none measures whether it still
\emph{knows} whether it succeeded. That gap is the subject of this paper, and
the answer is worse than a weak signal: the channel fails every acceptance
criterion we can import, every standard remedy for an unreliable channel fails
with it, and two unrelated model families are confidently wrong about the same
programs -- so a second opinion is not independent.

The obvious version of this measurement is a trap, and we avoid it. Stated
confidence on its own is a generated string, and there is direct evidence
against trusting it at this scale: a 0.5B self-judge has been reported to
perform no better than random while concentrating 60\% of its choices on a
single index~\cite{iscan_scaffoldnotvocabulary_2026}. We therefore do not read a
number out of prose. The readout is the normalized probability mass the model
places on \textsc{yes} against \textsc{no}, asked in both polarities and scored
against the realized hidden-test outcome. What that readout supports turns out to
be narrow, and we report it as such: it ranks correct work above incorrect, while
its calibration is not resolvable at this sample size.

\paragraph{Contributions.} All four are measurements of an instrument, and none
depends on an intervention producing a significant effect. That was deliberate: a
study whose contribution requires its treatment to work reports nothing when it
does not.

\begin{enumerate}
    \item \textbf{The first gauge study of an agent self-report channel, and all
        four elicitation methods in common use fail it.} On 120 frozen
        execution-graded programs judged through six semantically identical
        wordings, \%GRR runs $37.9$--$61.7\%$ against a $30\%$ rejection threshold
        and ndc $1.80$--$3.45$ against a required five. The same wordings on code
        the model did \emph{not} write fail by an indistinguishable margin, so this
        does not look like a fact about introspection
        (Section~\ref{sec:resolving}).
    \item \textbf{A planted effect whose measured size moves several-fold with
        wording alone}, $3.83\times$ on 37 execution-verified pairs and
        $4.17\times$ on an independent 46-pair sample, with a failed
        assay-sensitivity check: the effect ($+0.1962$) falls below the same
        instrument's own $SDC_{95}$ ($0.2776$) on the same parts, and no individual
        wording resolves it (Section~\ref{sec:surfacefloor}).
    \item \textbf{What that costs, and what fixing it costs.} Routing the
        least-confident $20\%$ of work to a human, changing only the wording, moves
        the share of genuine failures reaching that human from $26.5\%$ to
        $41.2\%$. Averaging $k$ wordings buys acceptance at $k{=}3$ for a logit
        readout and $k{=}7$ for a verbalized one -- but whether that also fixes the
        routing decision is unresolved at this sample size, and we say so rather
        than recommend the fix on the index alone (Section~\ref{sec:resolving}).
    \item \textbf{A released harness, and a record of what the check cost us.}
        The elicitation module, the frozen population, every measurement with the
        script that produced it, and checks that fail the build. Four headline
        numbers of our own were retracted by tests we ran against ourselves, and
        preparing this paper produced repeated instances of one defect -- comparing
        quantities whose measurements differed -- most found only once a guard for
        a previous one existed (Appendix~\ref{sec:repro}).
\end{enumerate}

Refusing to report an effect smaller than a measured floor is, in the four
fields named above, the ordinary rule rather than a strict one; for an accredited
laboratory it is mandatory and must be documented before any result is
reported~\citep{iso17025_2017}. Our contribution is that the check has an answer
in this domain, and what the answer is.

\paragraph{What we do not claim.} Not priority on placebo control for code
self-repair; that belongs to prior work cited above. Not that the placebo is
inert -- it cannot be proven inert, and Section~\ref{sec:interpretation} fixes
the exact sentences the design licenses. Not anything about frontier models,
multi-turn trajectories, tool use, repository-scale tasks, or cross-episode
memory: the tasks here are single-turn function repair on one frozen roster
with one frozen checkpoint. Not anything about consciousness, sentience, or
machine experience.

% ---------------------------------------------------------------------
\section{Related work}
\label{sec:related}
% ---------------------------------------------------------------------

\paragraph{Placebo control and nonspecific gains.} Prior art for the intervention
study we do not run, recorded because the apparatus is released and priority is not
ours. Three preregistered studies decompose self-repair feedback in frozen small
code models on HumanEval+ and MBPP+ with length-matched, derangement-assigned
placebos~\cite{iscan_scaffoldnotvocabulary_2026, iscan_falsificationnotexposure_2026,
iscan_formnotcontent_2026}; they own the task-success version of this experiment.
Corrupting a critique's correctness~\cite{stechly_gpt4doesntknow_2023}, randomizing
in-context labels~\cite{min_rethinkingdemonstrations_2022}, unrelated execution
feedback~\cite{gehring_rlef_2024}, irrelevant retrieved
documents~\cite{cuconasu_powerofnoise_2024}, token-matched vanilla
agents~\cite{hajimiri_onlineskillmodules_2026} and random-content retrieval
arms~\cite{wang_memdelta_2026} point the same way across adjacent subfields,
consistently enough that we treat an apparatus effect as expected rather than
surprising.

\paragraph{Elicitation and perturbation sensitivity.} Both artifacts we measure
have close neighbours, and the contribution lives entirely in the distinction, so
we draw it precisely. \citet{protocolsensitivity_2026} vary four measurement-protocol
axes with the elicitation prompt held fixed and report that the shift they induce
is comparable to the shift between model families -- noise as large as signal, in
numeric form. They then decline the next step: their limitations call the shifts
``sensitivity measurements, not estimates of how much each choice causally
contributes,'' and their ethics section disclaims deployable thresholds outright.
\citet{rescalingconfidence_2026} vary the response \emph{scale} rather than the
wording, and find verbalized confidence collapsed onto a handful of round numbers.
\citet{surfacevssemantic_2026} measure semantic against surface perturbation of
agent inputs at scale, and come nearest of anyone to our move: they subtract a
surface-perturbation inconsistency rate from the semantic one. Finally, averaging
a score across both presentation orders is already routine against ordering bias;
our polarity averaging is that correction on a new axis and we claim nothing for
it.

Two distinctions matter. First, \textbf{netting a noise level out of a point
estimate is not the same act as refusing to claim an effect that falls inside it.}
Subtraction yields a smaller number and still reports it; a gate can return
\emph{no claim}. Second, \textbf{none of this work constrains what may be
concluded}: we searched the full texts and found no admissibility rule, no
precondition, no effect withheld for being smaller than a measured floor. The gap
is not rhetorical -- \citet{rescalingconfidence_2026} report a scale effect whose
magnitude equals their own separately measured prompt-variation movement, and
comparing the two is simply not a step the field takes.


\paragraph{Floors already at work in agent evaluation.} Two 2026 papers put a
measured floor to use in this domain and are the closest prior art to what we do.
\citet{kaliyev_noisefloor_2026} measure a paired noise floor from a null
condition on multi-agent benchmarks and rule reported gains inadmissible against
it. \citet{qian_informationlimits_2026} preregister a three-valued verdict in
which effects inside a measured noise bar are declared not formally resolvable,
and hold to it against their own hypothesis. We do not compete with either for
the term; the difference is where the floor comes from. Both derive theirs from
\emph{sampling} variance -- seeds, replicates, runs -- which shrinks as
$\sqrt{n}$ and which a test-retest design can see. Ours comes from a
\emph{deterministic elicitation axis}: at temperature zero, re-running a phrasing
reproduces its answer exactly, so no replicate design can observe it and no
amount of additional data reduces it. That is a different variance component, and
it is the one a paired contrast cannot average away.

% ---------------------------------------------------------------------
\section{Measuring the readout before measuring the intervention}
\label{sec:floors}
% ---------------------------------------------------------------------

Every result below this line is about the instrument, not about reflection. All
of it is measured on frozen programs with frozen ground truth, so the code and
the correct answer are held constant and only the question varies. It is
reported first because it bounds everything after it.

\subsection{Elicitation}

The model is asked a binary question and the answer is read as normalized
probability mass on YES against NO over a single greedy token, rather than
parsed from a stated integer. Two details are not optional.

Mass must be summed over a \emph{token set}. The serving layer spreads the
answer across case and whitespace variants -- on one live call, \texttt{YES} at
0.977, \texttt{Yes} at 0.012, \texttt{" YES"} at 0.0007, against \texttt{NO} at
0.0093 -- so reading the single exact token reports a different number from the
one the model expressed. We sum over an allowlisted set, normalize over the
decisive mass, and report what fell outside it. Across the 40 readouts whose
per-call mass we persisted, the probability landing on neither answer -- an
upper bound on what escaped the requested top-$k$ -- stays below
$6\times10^{-4}$.

A readout with no mass on either answer is a refusal, not a coin flip. The
implementation raises rather than returning $0.5$, because a substituted $0.5$
is indistinguishable in the analysis from a genuine one.

\subsection{Polarity bias}

Asking "will this PASS?" and asking "will this FAIL?" are the same question. On
identical programs, with everything else held fixed, the two answers do not
agree. Averaging a polarity and its complement,
\begin{equation}
    p = \tfrac{1}{2}\left[\,P(\textsc{yes}\mid\text{will it pass})
        + \bigl(1 - P(\textsc{yes}\mid\text{will it fail})\bigr)\right],
\end{equation}
moves the readout by a mean of \textbf{0.103} against the single-polarity
version. That displacement is a property of the question. It also improves the
readout on every axis we measure: separation between correct and incorrect
programs rises from $+0.040$ to $+0.0705$, AUROC from 0.693 to 0.713, and Brier
falls from 0.2455 to 0.1930. These are the $n{=}40$ self-assessment
sample; the larger replication of Section~\ref{sec:floors} reports the same
readout at $n{=}120$.

The consequence for prior work is direct. A single-polarity probe carries this
bias into every condition it measures, and 0.103 is larger than the effects such
studies report.

\subsection{The phrasing floor}

Polarity is one axis of variation. The general version asks how far the readout
moves under \emph{any} semantically empty change to the question. We hold the
programs and the ground truth frozen and vary only the wording,
across six phrasings that ask the identical question with the identical
answer space -- including the one this paper otherwise uses, so our own
configuration sits inside its own floor rather than being privileged.

Over 120 frozen, execution-graded programs spanning both benchmarks, with the
same model, seed and decoding, the six phrasings return AUROCs from 0.6091 to
0.7370 (spread 0.1279) and Brier scores from 0.1878 to 0.2342 (spread 0.0464). On
an independent 40-program sample the same six return 0.6967 to 0.8067 (spread
0.1100). We report the second sample because the two do not agree, and we present
these as text rather than as a table because -- as the next three paragraphs
establish -- the spread is not a measured effect and the ranking does not
reproduce. Per-phrasing values are in the released artifact.

These figures are where an earlier version of this paper went wrong, and the
correction is instructive enough to keep in the body.

\textbf{The aggregate spread is not significant, at either sample size.} A
within-item permutation null -- shuffling phrasing labels within each program,
which preserves every program's score profile and destroys only the alignment of
phrasings across programs -- gives, at $n{=}120$, a null median of 0.0954 against
an observed 0.1279, $p = 0.178$. Per benchmark: HumanEval+ observed 0.1457
against null 0.1538 ($p = 0.555$), MBPP+ observed 0.0806 against null 0.1258
($p = 0.883$). At $n{=}40$ the observed spread sat \emph{below} its null median
entirely. We previously reported the spread as a measured phrasing effect. It is
not one.

\textbf{The per-phrasing ranking is a draw, and that is the sharper evidence.}
Spearman $\rho$ between the two samples' orderings is $+0.029$. The
\texttt{code-first} phrasing is best at $n{=}40$ (0.8067) and worst at $n{=}120$
(0.6091). A ranking that does not reproduce between two samples of the same
population is not a property of the phrasings, and a paper reporting one such
ordering -- ours included -- reports a draw from a distribution it does not
characterise.

\textbf{Can you tune the wording on a held-out set?} Not measurably: choosing
the best-discriminating wording on 60 held-out problems identifies the best on
the remaining 60 no more often than guessing, and the interval on that is wide
enough that the study cannot tell whether tuning helps at all
(Appendix~\ref{sec:extra}).

\textbf{What does survive gets stronger with data.} Pooled discrimination is
AUROC \textbf{0.7172} with a bootstrap interval of \textbf{[0.608, 0.819]},
excluding chance and tighter than at $n{=}40$. The Brier gap against a constant
predictor is $+0.0029$, interval $[-0.023, +0.030]$: at triple the sample,
calibration remains \emph{unmeasured} rather than shown to be poor.

\textbf{The reconciliation is the finding.} Per-item, the readout is unstable --
standard deviation 0.0497 across phrasings, maximum range 0.5207 on identical
programs. Aggregate, it is steady. Both are true because aggregate AUROC averages
over items, and the variance decomposition puts the dominant component in the
item\,$\times$\,phrasing interaction (0.0046) rather than the phrasing main
effect (0.0018). Averaging over items averages the interaction away.
\textbf{A paired contrast does not.} So a study reporting aggregate
discrimination on this channel would see nothing wrong with it, while a study
reporting a paired intervention effect on the same channel is exposed --
which is why the gauge indices of Section~\ref{sec:resolving}, and not this
table, carry the paper.

\subsection{Discrimination without calibration}
\label{sec:dissociation}

On its own graded output -- the model writes a program, the EvalPlus harness
grades it, and the model is then asked about that program -- the readout
separates correct from incorrect at \textbf{AUROC 0.747} on 40 programs, with
a bootstrap interval of \textbf{[0.555, 0.913]} that excludes chance; the
$n{=}120$ replication of Section~\ref{sec:floors} returns 0.717 with a
tighter interval. Its calibration is a
different matter: the Brier gap against a constant-base-rate predictor is
$+0.0066$ with an interval of $[-0.029, +0.046]$ and $P(\text{worse}) = 0.62$.
\textbf{We therefore claim that the readout ranks, and we claim nothing about
whether it estimates} -- at this sample size the calibration question is simply
unanswered, and an earlier draft of this paper overstated it as a demonstrated
failure.

Individual cases show what the ranking is made of. A program passing
\textbf{0 of 105} hidden tests drew $P(\textsc{yes}) = 0.922$; one passing 9 of
110 drew 0.969. In the other direction, two programs that passed every test drew
0.016 and 0.037. These are illustrative rather than inferential.

The operational consequence follows from the part that is established. A pipeline
that \emph{ranks} by agent-stated confidence -- review the least confident decile
-- is consuming a signal that demonstrably exists. A pipeline that
\emph{thresholds} on it -- escalate below 0.8, proceed above -- depends on
calibration these data cannot support, and on an absolute scale whose position
moves with the wording of the question. Those two uses are routinely treated as
one capability, and only the first is warranted here. Binning by stated confidence shows
both halves at once. Accuracy does rise with confidence in both families
(Appendix~\ref{sec:extra}), so ordering by confidence orders by risk within a
model. \textbf{But the two scales do not overlap at all.} One
family's bins average $0.688$ to $0.929$ and the other's $0.289$ to $0.523$:
the \emph{least} confident programs of the first score higher than the
\emph{most} confident of the second. Any threshold between those ranges marks
every program from one model safe and every program from the other risky,
whatever their correctness.

\textbf{What the warranted use costs.} Ranking survives, so we priced it. Send
the least-confident $20\%$ of programs to a human and change nothing but the
wording of the question: the share of genuine failures that reaches that human
moves from $26.5\%$ to $41.2\%$ -- \textbf{one wording delivers half again as
many real defects to review as another, at identical cost} -- and the two
wordings agree on only half the programs they send. The spread is $14.7$ points
with a bootstrap interval of $[6.1, 33.3]$ over programs, excluding zero, and it
holds at $10\%$ and $30\%$ budgets too. This is the gauge failure denominated in
the currency a verification pipeline actually spends, and the choice driving it
is one no reporting convention asks anyone to write down.

\subsection{One planted effect, measured six ways}
\label{sec:surfacefloor}

The results so far concern the readout in isolation. The question that decides
whether any of it matters is different: does the measured \emph{size of a real
effect} depend on how the question is worded? A phrasing main effect would not
matter, because it is common-mode and cancels exactly in a paired contrast. An
interaction with the arm would matter, because it does not.

We plant an effect rather than hoping to find one. Each problem contributes two
variants -- the canonical solution and a deterministically mutated version -- and
\textbf{both are graded by the execution harness}, so a pair enters the study only
when the canonical passes every hidden test and the mutant genuinely fails one.
That verification matters, and we ran it on two independently drawn samples: 4 of
50 candidate mutants still passed in one (8.0\%) and 3 of 40 in the other
(7.5\%). Three of the four are comparison-boundary rewrites, where $\ge$ to
$>$ changed nothing because the hidden tests never exercise the boundary; the
fourth narrows a loop range and is a different failure, which is why we report
the rule alongside every dropped pair rather than a single explanation. Those pairs
would have contrasted programs that behave identically. We report the 46 and 37
survivors respectively: the larger sample gives the wording spread, and the
smaller one additionally carries its own $SDC_{95}$ over the same parts, which is
what an assay-sensitivity comparison requires and the larger sample cannot supply.
Each variant is judged under all six phrasings, and the estimand is the paired
difference
$d = \overline{P(\textsc{yes}\mid\text{correct}) - P(\textsc{yes}\mid\text{mutated})}$,
which is exactly the shape of quantity a placebo-controlled study reports.

\begin{figure}[t]
    \centering
    \includegraphics[width=0.86\linewidth]{fig_planted_effect.pdf}
    \caption{One planted effect, measured six ways, over the 37
    execution-verified pairs that also carry their own floor. Each point is the
    paired difference $d$ between the correct and
    mutated variants of the same programs under one wording; bars are standard
    errors over problems. Programs, ground truth, model, seed and decoding are
    identical across rows -- only the wording differs. The shaded band is
    $SDC_{95}$ measured on \emph{these} parts through \emph{these} wordings, so
    points inside it are effects this instrument cannot resolve. \textbf{All six
    are inside it.} An earlier version of this figure drew the band from a floor
    measured on a different task under a different question wording, against which
    four of the six appeared to clear it.}
    \label{fig:planted}
\end{figure}

\textbf{How big is the wording effect next to the choice everyone already
controls for?} Smaller, and we report that. Two wordings of one model agree
across the 120 programs at $r = 0.863$; two model families on the \emph{same}
wording agree at $0.536$, a difference of $0.33$ with interval
$[0.25, 0.40]$. The model is decisively the larger source of variation.
\textbf{The asymmetry is in what gets written down.} Every paper reports which
model it used; none reports the wording. A smaller source of variation that no
convention controls for is the worse problem, not the lesser one.

\textbf{Figure~\ref{fig:planted} is the paper's central claim about
consequences.} The identical effect measures $3.83\times$ larger under one
wording than another, against standard errors an order of magnitude too small to
account for it. Two laboratories running the same experiment on the same
programs and reporting the same estimand would differ several-fold for no reason
but how they asked the question -- and neither would report the phrasing, because
no convention requires it. The failure does not depend on averaging: across the
six wordings the effect runs from $+0.0643$ to $+0.2465$ and \emph{every one}
falls below the $SDC_{95}$ of $0.2776$ measured on those same verified parts, so no
choice of question would have resolved this effect.

\textbf{How precisely is that ratio known?} Not to three significant figures. A
max/min over six wordings takes both ends from extreme order statistics and
compounds their noise -- exactly the statistic a paper about overstated
resolution should distrust in its own results. Resampling items and recomputing
both extremes per draw gives a 95\% CI of $[2.71\times, 7.51\times]$ here, and
$[3.09\times, 6.83\times]$ on the independently drawn 46-pair sample that supplies
the wording-spread figure quoted in the abstract. Neither interval is narrow, both
exclude $1\times$ by a wide margin, and they overlap almost entirely -- which is
the claim that matters: the distortion is multiplicative, at least threefold, and
it replicates across samples. We report both rather than the larger.

\textbf{Assay sensitivity fails on a matched comparison.} A
smallest-detectable-change is only a comparator for an effect measured by
\emph{the same instrument} on \emph{the same parts}, so we report the two together
from one run. Over 37 canonical/mutant pairs whose mutants were confirmed broken by
the execution harness, read out through six wordings, qwen2.5-coder:7b returns
$SDC_{95} = 0.2776$ against a planted effect of $+0.1962$: the instrument cannot
resolve a content difference that execution establishes by construction. The
second model family agrees, at a tenth the effect size
($+0.0188$ against $SDC_{95} = 0.1171$), though we hold that corroboration to a
weaker standard and say so: its mutants were not put through the execution harness,
and grading the qwen pairs the same way removed 3 of 40 as still passing. By the standard logic of ICH
E10~\citep{ich_e10_2000} a null from such an instrument is
\emph{uninterpretable} -- it does not separate ``no effect'' from ``no
sensitivity'' -- and that verdict binds the intervention study this harness was
built to run exactly as it binds anyone else's.

\textbf{We arrived at this comparison by failing to make it correctly first.} An earlier draft judged this effect against a floor measured on a different task, invisible to every review and to a traceability audit that passed at $109/109$; the borrowed floor was the lower of the two, so the error favoured us. Seven instances of that one defect surfaced while preparing this paper, and the checks that now fail the build on it are in Appendix~\ref{sec:repro}.

\textbf{Verification moved the numbers in the direction it had to, and changed no conclusion} (Appendix~\ref{sec:extra}).

\textbf{The interaction is real, and we report its size honestly:} phrasing is $8.7\%$ of paired-effect variance, a modest share of a total dominated by between-problem heterogeneity yet large against the effects anyone is trying to estimate (Appendix~\ref{sec:extra}).

\subsection{Which elicitation method resolves anything?}
\label{sec:resolving}

The verdict above is about one readout. The field uses several, and papers choose
between them on convenience rather than on resolving power, because resolving
power has not been reported for any of them. Everything in this subsection is
measured on the \textbf{self-assessment} task -- the model judging code it wrote
itself, with correctness established by the execution harness -- because that is
the task a verification pipeline actually consumes. We therefore run the identical gauge
study over four: the verbalized integer that dominates the literature, a
verbalized five-level ordinal, and the single-polarity and symmetrized logit
readouts. Same parts, same appraisers, same programs, same ground truth; only the
question format changes.

\begin{table}[t]
    \centering
    \caption{Gauge study over four elicitation methods, on the
    \textbf{self-assessment} task: the model judges programs it generated itself,
    graded by the execution harness. 120 execution-graded programs as parts, six
    semantically identical phrasings as appraisers, one trial per cell. AIAG
    rejects a measurement system above 30\% GRR and requires at least five
    distinct resolvable categories. \textbf{No method passes either criterion.}
    These figures are not comparable with a gauge study run on other-code
    judging, which is the easier task (AUROC 0.776 against 0.734 on this
    checkpoint, both scored the same way over the same six wordings). Bold marks the best value in each column. All four rows come from
    one run under one appraiser set; an earlier 40-part version of this table
    stitched the symmetrized row in from a separate measurement after we found
    our pipeline had filled it by reusing an artifact produced with a different
    appraiser set, and no verdict changed at either size.}
    \label{tab:resolving}
    \begin{tabular}{lrrrrrr}
        \toprule
        Elicitation method & ICC & SEM & $SDC_{95}$ & \%GRR & ndc & AUROC \\
        \midrule
        Verbalized integer 0--100   & 0.657 & 0.0693 & 0.1922 & 58.5 & 1.95 & 0.670 \\
        Verbalized ordinal (5-level)& 0.619 & 0.0989 & 0.2741 & 61.7 & 1.80 & 0.656 \\
        Logit, single polarity      & \textbf{0.857} & 0.0596 & 0.1652 & \textbf{37.9} & \textbf{3.45} & 0.726 \\
        Logit, symmetrized          & 0.819 & \textbf{0.0527} & \textbf{0.1461} & 42.5 & 3.00 & \textbf{0.734} \\
        \midrule
        \multicolumn{4}{l}{\emph{AIAG acceptance}} & $<30$ & $\ge 5$ & \\
        \bottomrule
    \end{tabular}
\end{table}

Table~\ref{tab:resolving} is the paper's central result. Three things follow.

\textbf{Nothing passes.} Every method is rejected on both criteria, and not
narrowly: the best resolves about 3.5 distinct categories of the quantity it is
supposed to measure, the worst 1.80, against a required five. Parse failures were
zero throughout, so this is not an artifact of malformed output.

\textbf{And it is not, so far as we can tell, a fact about introspection.}
Running the same six wordings over code the model did \emph{not} write --
canonical against mutated programs, correctness fixed by execution -- returns
the same verdict at the same magnitude: \%GRR $42.9$ against $42.5$, ndc $2.97$
against $3.00$, $ICC$ $0.816$ against $0.819$. Every difference is
indistinguishable from zero. We are careful about what that licenses: the
intervals are wide -- the \%GRR difference is $[-8.7, +8.3]$ -- so a moderate
gap would go undetected, and this bounds how special self-report could be rather
than showing the two are identical. It is enough to say the failure does not
look like a property of self-knowledge, and that a verifier consuming any binary
LLM judgement about code inherits it.

\textbf{The channel the field reports is the weaker one.} At this sample the
verbalized integer is worse than both logit readouts on all six axes -- \%GRR
58.5 against 37.9 and 42.5, ndc 1.95 against 3.45 and 3.00, $SDC_{95}$ 0.192
against 0.165 and 0.146, AUROC 0.670 against 0.726 and 0.734 -- and the ordinal
variant is worse still, at 61.7\% GRR. We flag that this was not true at
$n{=}40$, where the integer's compressed scale gave it the smaller absolute
error while it lost everywhere else; verbalized confidence collapses onto a
handful of round values~\citep{rescalingconfidence_2026}, and compressing a
scale shrinks absolute spread and the capacity to separate parts together. A
study that elicits confidence as a number in prose is choosing the least
resolving option available to it, and nothing in the literature would have told
it so. Nor does prose capture what the logits miss: the channels are far from
redundant yet combining them adds nothing measurable
(Appendix~\ref{sec:extra}).

\textbf{The two logit readouts split by criterion, and the split is not stable.}
Single-polarity takes the ratio indices (\%GRR 37.9 against 42.5, ndc 3.45
against 3.00); symmetrized takes absolute resolution ($SEM$ 0.0527 against
0.0596, $SDC_{95}$ 0.146 against 0.165) and, at this sample, discrimination
(AUROC 0.734 against 0.726). The disagreement is structural rather than
contradictory: symmetrizing removes a polarity bias that inflates between-part
spread, and the ratio indices divide by exactly that spread, so a method can
shrink its own measurement error and still lose ground on \%GRR. The AUROC
ordering reversed between $n{=}40$ and $n{=}120$, which is worth reporting rather
than smoothing: a ranking that moves between samples of the same population is
the thing this paper argues nobody should treat as a stable property.

\textbf{What would it take to pass?} The failure is driven by variance across
wordings, so averaging $k$ wordings per judgement divides measurement variance
by $k$ while leaving between-part variance alone, and both indices improve. The
smallest $k$ clearing both is $3$ for either logit readout against $7$ and $8$
for the verbalized integer and ordinal. Since $k$ counts \emph{wordings} and the
symmetrized readout asks both polarities of each, the cost in model calls is $3$,
$6$, $7$ and $8$ respectively. \textbf{Single-polarity logit reaches acceptance
for three calls per judgement; the verbalized integer needs seven.} Symmetrizing
doubles the per-wording cost to land within one call of the prose readout it is
meant to dominate, and it already trailed single-polarity on both ratio indices,
so the cheap route to acceptance is the single-polarity readout. Held at a fixed call budget the same answer appears, and buying wordings
rather than polarities is what improves the indices (Appendix~\ref{sec:extra}).

\textbf{A second model family: the verdict replicates, the price does not.}
Running the identical gauge over llama3.1:8b rejects all four methods again
(\%GRR $60.4$--$93.5$, ndc $0.54$--$1.86$) and preserves the ordering, so
preferring logits is not an artifact of one model. The design cost is another
matter: $88$ calls per judgement for the verbalized integer against $7$ for the
same readout on qwen, $28$ against $8$ for the ordinal, $8$ against $3$ for
single-polarity logit. \textbf{Between $2.7\times$ and $12.6\times$ apart for
the same method, so $k$ cannot be read out of this paper and applied} -- it is a
property of a deployment and has to be measured there, which makes the design
cost a procedure rather than a recommendation. On the format this literature
reports most often, that second model returns AUROC $0.514$ and $ICC$ $0.126$:
ranking at chance, and a readout that would need $88$ calls to become
admissible.

\textbf{And the two families are wrong about the same programs.} The
standard remedy for an unreliable verifier is a second opinion -- run two
models, escalate on disagreement -- which assumes the errors are independent.
They are not. Restricting to the 34 failures, where every program is equally
wrong so the label cannot manufacture agreement, the two families' readouts
still correlate at $+0.63$: they agree on which failures \emph{look} correct.
Of each family's eleven most-overconfident failures, eight are the same
programs, against $3.56$ expected by chance ($p = 0.001$; Jaccard $0.571$,
bootstrap interval $[0.333, 0.778]$, excluding the chance value $0.193$).
\textbf{On this task a second model is not an independent second opinion.} Agreement reads as
corroboration and is exactly where both models are confidently wrong, which
leaves disagreement-based escalation blind precisely where it is needed.

\textbf{Every standard remedy we could test fails, and they fail
differently.} A second model does not help, because the families are
overconfident about the same programs. A threshold does not help, because the
two scales do not overlap. \textbf{Calibration does not help either}: a Platt
or isotonic map fitted on half the programs and applied to the other half does
not beat predicting the base rate, in either family (base rate minus
calibrated Brier, $+0.022$ with interval $[-0.014, +0.037]$ and $+0.005$ with
$[-0.021, +0.014]$). For the second family the uncalibrated readout is worse
than a constant outright, Brier $0.296$ against $0.208$ -- using that
confidence as a probability is worse than ignoring it. Averaging wordings is
the one remedy that moves an acceptance index, and its effect on the routing
decision is unresolved. What survives is ranking within a model, which we now
report with the alternatives eliminated rather than merely unexamined.

\textbf{And the remedy that works does not reach the failure that matters.}
Averaging wordings repairs the acceptance indices, so it is natural to hope it
also disperses the shared blind spot -- if which failures look correct depended
on the wording, averaging would scramble it. It does not. Across every subset
of every size, the cross-family overlap does not fall as wordings are averaged;
it runs from $6.5$ of $11$ at $k{=}1$ to $8$ at $k{=}6$, against a chance level
of $3.6$, and the overlap after averaging all six remains above chance
($[+1.2, +3.9]$). We do not claim averaging makes it \emph{worse}: the
increase itself is not established ($+0.70$, interval $[-0.33, +1.83]$).
\textbf{The claim is that averaging does not disperse it}, which is enough --
the one remedy that repairs the measurement leaves the danger exactly where it
was. A mechanism would explain the direction, since wording noise is the only
source of disagreement between the families about which failures look correct,
so averaging it away can only leave the shared component; we note it as a
possibility rather than a result. The obvious
explanation is the near-miss -- a program failing one test of a hundred looks
correct -- and grading every failure against its hidden tests rules it out. The
eight shared blind spots pass a mean $0.317$ of their tests against $0.471$ for
the failures both families correctly doubted, a difference in the wrong
direction for that story and not significant either ($p = 0.34$). What remains
is descriptive, needs no test, and is stronger than the mean suggests. Not one
of the eight fails a single test -- the profile a flaky or over-strict
assertion would produce, and the reading that would make the models right and
the harness wrong. Seven of eight fail more than three. \textbf{Three of them
pass zero tests out of $112$, $116$ and $150$, and a fourth passes one of
$106$.} Two architectures trained on different data both rank programs that
pass \emph{none} of their tests among the most likely to pass all of them, and
rank the same ones. We verified this rather than extrapolating it: building the
averaged design from every disjoint partition of our six wordings reproduces the
predicted medians. \textbf{And it carries the paper's own warning.} Whether
$k{=}3$ actually passes depends on which three wordings were grouped -- \%GRR
spans $[17.6, 37.4]$ across the ten partitions, straddling the threshold, with
only $60\%$ of groupings clearing it. The remedy inherits the instability it is
meant to remove, and a single reported $k$ would hide that.

\textbf{And the index is not the decision.} At $k{=}3$ the gauge indices cross
into acceptance, so the natural reading is that three calls solve the problem.
Repeating the routing analysis over averaged readouts does not support that: the
spread in genuine-failure recall falls from $14.7$ points to $8.8$, but paired
within bootstrap worlds that reduction is $+0.039$ with a $95\%$ interval of
$[-0.083, +0.182]$ and $P(\text{reduction}) = 0.65$. \textbf{Averaging fixes the
index; whether it fixes the decision is unresolved at this sample size.} We note
that comparing the two marginal intervals would have misled in both directions --
they overlap almost entirely, inviting \emph{no difference}, while the point
estimates invite \emph{$40\%$ better} -- and only the paired difference answers
it, both statistics being computed on the same programs. This is the paper's own
finding one level up: the study that would tell you whether the remedy works is
underpowered in the same way as the study it is meant to rescue.

% ---------------------------------------------------------------------
\section{The study this instrument cannot support}
\label{sec:population}
\label{sec:arms}
\label{sec:placebo}
\label{sec:interpretation}
\label{sec:placebo3}
% ---------------------------------------------------------------------

This work began as a placebo-controlled decomposition of post-failure
reflection. The apparatus is built, tested and released; reporting what happened to
it is more useful than reporting the study.

\textbf{What was built.} A seven-arm, parity-asserted design over EvalPlus repair
tasks~\citep{liu_evalplus_2023}, built, tested and released; the arm roster, the
parity classes and the frozen pool are in Appendix~\ref{sec:design}.

\textbf{Why it is not the paper.} The estimand is a paired difference between
arms on the self-report channel, and Section~\ref{sec:surfacefloor} shows the
instrument does not resolve an effect planted by construction on the easier task.
Published reflection effects are smaller still. A null here would not distinguish
``the placebo and the real reflection do the same thing'' from ``the instrument
cannot tell them apart'' -- a lack of assay sensitivity~\citep{ich_e10_2000} --
and we record that before running the arms rather than afterwards to explain a
null.

\textbf{What would have to change first.} Three things, in order of how much they
buy. A readout whose $SDC_{95}$ is smaller than the effect of interest -- none of
the four we measured qualifies, so this means a new elicitation method rather
than a different question. A positive control that passes, since without one no
null is interpretable. And an estimand chosen with the interaction in mind:
because the phrasing main effect cancels in a paired contrast while the
arm\,$\times$\,phrasing interaction does not, a design that randomises phrasing
across items and reports the interaction is strictly better than one that fixes a
single wording and does not mention it.

\textbf{What is released.} The arm builders and their parity asserters, the
frozen pool and its digest, the donor machinery and leakage audit, the
subprocess-isolated grader, the elicitation module, and every measurement in this
paper with the script that produced it. A study wanting to run this design does
not have to rebuild it -- but should first measure whether its own readout can
resolve what it intends to claim.

\section{Limitations}
% ---------------------------------------------------------------------

\textbf{One model.} Every number is Qwen2.5-Coder-7B-Instruct~\citep{hui_qwen25coder_2024} at Q4\_K\_M under
greedy decoding. A gauge verdict is a property of a measurement system, and this
is one measurement system. We report a cross-family replication where available
and treat the single-model result as a hypothesis about the channel rather than
a fact about it. Quantization is unexamined: all four elicitation methods were
measured at the same 4-bit precision, and whether the readout's dispersion is a
property of the model or of its quantization is untested.

\textbf{Sample size, and what it does and does not bound.} The gauge study uses
40 parts and the floor replication 120. Variance components are estimable at that
size; the aggregate significance tests are not powerful, and we say so where it
matters -- the phrasing spread is reported as \emph{not demonstrated} rather than
as absent, and the calibration gap as \emph{unmeasured} rather than as zero.

\textbf{One trial per cell.} Repeating a probe shows only the \emph{first} call for a prompt differs; the rest are bit-identical and the seed is inert, so \%GRR is reproducibility plus interaction with no repeatability term. Call order contributes $0.8\%$ of the appraiser variance (Appendix~\ref{sec:extra}).

\textbf{Six phrasings, authored not sampled.} The appraiser facet is a convenience set written by us, not a draw from any defined population of reasonable wordings, so it estimates phrasing dispersion and does not bound it. We characterise it rather than only disclaiming it: mean pairwise content-word overlap is $0.18$ (range $0.06$--$0.43$) across a $2.55$-fold length ratio, so these are substantially different wordings rather than near-clones. That reading is the one against us. Near-clones would have made the measured spread a lower bound with the true exposure worse than shown; a diverse set instead models what independent authors at different labs would write, and a single lab nudging its own wording should expect less.

\textbf{Mutant validity is measured, not assumed.} Both variants of every pair
are graded by the execution harness, and a pair is used only when the canonical
passes and the mutant genuinely fails. Across two independently drawn samples
4 of 50 candidate mutants still passed and 3 of 40 did, 8.0\% and 7.5\% --
close enough to read as a property of the mutation rules rather than of one
draw. The residual risk is the reverse case -- a mutant that fails
for a reason unrelated to the planted change -- which grading cannot distinguish
and which we do not claim to exclude.

\textbf{The blind spot is demonstrated for self-assessed code, and only partly
beyond it.} The 120 programs were generated by the first model, so one family is
assessing its own output while the other judges code it did not write. We ran the
clean version on 37 execution-verified canonical/mutant pairs, which neither
model produced. Half of it replicates: on mutants alone, where every program is
equally broken, the two families' readouts still correlate at $+0.475$ against
$+0.63$ on the self-assessed set. Half does not: the overlap of most-convincing
mutants is $5$ of $12$ against $3.9$ expected, $p = 0.32$. \textbf{So the sharp
form of the result belongs to self-assessment}; on code neither model wrote the
agreement survives as a correlation while the tail overlap is unresolved at this
sample size.

\textbf{The cross-model comparison is the easier task.} To hold part variance
constant across models, every model judges the same fixed programs rather than its
own output. Our own measurements put other-code judging at AUROC 0.776 against
0.734 for self-assessment on the same checkpoint, each scored as the mean readout
over the same six wordings. Failures measured that way are a
\emph{lower bound} on what the self-report channel would show.

\textbf{Two claims were retracted during this work and the paper is their record.}
An earlier draft reported the study-level AUROC spread across phrasings as a
measured effect, and reported the readout as worse-calibrated than a constant.
A within-item permutation test ($p = 0.178$ at $n{=}120$, with the $n{=}40$
observation below its own null median) and a bootstrap interval spanning zero
retired both. We state this in the body rather than silently correcting it,
because a paper arguing that instruments overstate has no standing to quietly
restate its own numbers.

\textbf{What the gauge verdict does not establish.} That a channel fails \%GRR and
ndc means it cannot resolve differences of the size those indices imply. It does
not mean the underlying quantity is absent, that no elicitation method could
resolve it, or that a larger model would fail identically. It is a statement about
what can be concluded from this instrument, which is the only kind of statement a
gauge study makes.

\textbf{The intervention study is not run.} The placebo decomposition described in
Sections~\ref{sec:arms}--\ref{sec:placebo} is specified and its harness is built
and tested, but no arm has been executed. Every quantity attached to it is a
placeholder. Given $SDC_{95} = 0.22$, we expect it to return \textsc{unresolvable}
on the primary outcome, and that expectation is recorded before the fact rather
than offered afterwards as an explanation.

\textbf{Two implementation limits, recorded rather than papered over.} The grading
child runs candidate programs without a hardened sandbox, in a temporary working
directory with ordinary process privileges; acceptable for model-generated repairs
to these benchmarks, unsafe for adversarial input. The parent timeout backstop is
heuristic and could false-timeout a problem with many genuinely slow passing tests.

\section{Conclusion}
% ---------------------------------------------------------------------

A common answer to who verifies the agents is that the agent helps: it reports a
confidence, it asks for review when unsure. That answer assumes the report is a
window onto the agent's state rather than an output of the apparatus that elicited
it. We did not test whether the window is clean. We tested whether it is a window
at all, using procedures four fields already require and none had applied here.
This channel does not meet the entry criteria those fields set: it is rejected by
both industrial gauge indices in every elicitation format we tried; it cannot
resolve a content effect handed to it by construction; and the measured size of
that effect moves at least threefold with the wording of the question alone. The
last has the teeth, because wording is exactly the design choice no reporting
convention asks anyone to disclose.

None of this shows that agents lack self-knowledge, and we are careful not to say
so. There is mechanistic evidence the other way: verbal confidence is computed at
answer-adjacent positions and cached before it is verbalized, and those
representations explain variance beyond token
log-probabilities~\citep{kumaran_verbalconfidence_2026}. That is what makes an
unresolving instrument worth reporting rather than expected -- something is there,
and the instrument used to look for it cannot resolve differences of the size being
claimed, which makes the existing evidence -- ours included -- weaker than it
appears rather than wrong. The remedy is a better instrument, and a
habit of measuring it before trusting it.

We started this work intending to run a placebo-controlled intervention on this
channel. The instrument check said the result would be uninterpretable, so we
are reporting the check instead. That is the smaller paper, and it is the one the
evidence supports.

\bibliographystyle{plainnat}
\bibliography{refs}

\appendix
\section{Reproducibility statement}
\label{sec:repro}
% ---------------------------------------------------------------------

The failure pool is frozen and digest-bound before any arm exists, and nothing
downstream may add to it. Model selection ran only an initial attempt, a plain
retry, and an oracle-hint retry on a disjoint calibration slice; no real or
placebo reflection was generated during selection, and the selection rule
(largest eligible, never largest uplift) was fixed in advance. Donor assignment
is a published seeded derangement with a manifest digest that fails on any
post-hoc edit. Parity asserters return measured records rather than silently
repairing defects. No LLM judge appears in grading, in any outcome, or in any
validity check, and the constant recording this is asserted by test and
serialized into every record. The analysis script is run exactly once, after
the preregistration is frozen to a commit hash. Harness, pool, manifests, and
the \textsc{Placebo-3} implementation are released.

\paragraph{Checks released with the paper, and why there are five.}
Preparing this paper produced seven instances of a single defect: two quantities
compared as commensurable when the measurements behind them differed. They appeared
in the abstract, the contributions list, the body, the central gauge table, the
central figure, and twice in a limitation. Four were found only after a mechanical
guard for a previous one existed, which is the argument for having built them. All
five checks run in the pre-submission gate and each is stated with what it does not
catch, because a guard advertising coverage it lacks is how the earlier instances
survived review.

\begin{enumerate}[itemsep=1pt, topsep=3pt]
    \item \textbf{Prompt fingerprints.} Every elicited measurement carries a digest
        of the exact templates that produced it, and a comparison across differing
        digests raises rather than warns. Deliberately brittle: it does not
        normalise case or whitespace, because this paper's own result is that such
        differences move the readout, so a fingerprint that forgave them would
        certify the comparisons most likely to be unsound.
    \item \textbf{Cross-artifact reuse.} Any value repeated across artifacts to full
        floating-point precision is a copy, not a second measurement: two runs of
        one quantity agree to a few decimal places, never to seventeen. Legitimate
        sharing is declared with a reason and each declaration is narrow -- it does
        not excuse the same value appearing in a third artifact.
    \item \textbf{Floor bindings.} Each detection floor is bound to the artifact it
        must come from and the two are required to match exactly. A floor is the
        one quantity whose provenance changes a verdict, since an effect is
        resolvable or not depending on what it is judged against.
    \item \textbf{Unmatched comparisons.} If the text says $X$ falls below $Y$, some
        artifact should contain both; when none does, the sentence asserts a
        relationship nothing measured. Cross-model, cross-task and before/after
        comparisons are exempt by construction, and every exemption is a blind spot
        we name rather than hide.
    \item \textbf{Numeric traceability.} Every number in the built PDF must appear
        in a persisted artifact. This check is the weakest of the five and we say
        why: it passed on a draft in which an AUROC was attributed to a task it had
        never been measured on, because the value had been added to the allowlist of
        externally-sourced constants. No measured value may sit in that list, and
        the audit now runs with none.
\end{enumerate}
\TODO{anonymous artifact identifier and bundle hash}


\section{Additional measurements}
\label{sec:extra}
% ---------------------------------------------------------------------

Each of these answers a question a reader may reasonably ask of the main
result. None carries a step of the argument, which is why they are here.

\paragraph{Tuning the wording on a held-out set.} An unstable ranking licenses
someone to try exactly this, and it is where the result becomes a decision
rather than a correlation. Re-analysing the same 120
readouts: choosing the best-discriminating wording on 60 held-out problems
identifies the genuinely best wording on the remaining 60 in $15.5\%$ of splits,
against $16.7\%$ by guessing, and gains $+0.015$ AUROC over picking one at
random. Both numbers invite the conclusion that selection is worthless, and
\emph{neither survives its own interval}: resampling programs rather than splits
-- splits reuse one fixed program set and are not independent observations --
gives $[0.015, 0.800]$ for the hit rate and $[-0.010, +0.059]$ for the gain. The
correct statement is the weaker and more useful one: at this sample size the
study cannot tell whether tuning the wording helps at all -- the same resolution
failure the paper measures elsewhere, arriving at a decision someone makes on a
Tuesday.

\paragraph{Do the channels carry different information?} The two
families are far from redundant -- verbalized ordinal correlates with
single-polarity logit at $\rho = 0.16$, against $0.83$ between the two logit
readouts -- yet the best combination of channels beats the best single channel
by $+0.002$ AUROC, $95\%$ CI $[-0.052, +0.062]$. The prose channel carries
something different; it does not carry something that predicts correctness.

\paragraph{We arrived at this comparison by failing to make it correctly first, and
the failure is worth more than the fix.} An earlier draft judged this effect
against a floor measured on the self-assessment task under a different question
wording. Both were labelled with the same six phrasing names, so the mismatch was
invisible to every review we ran and to a traceability audit that passed at
$109/109$ -- a copied number traces perfectly, because it really is in an
artifact. The borrowed floor was the \emph{lower} of the two, so the error had
been working in our favour. Seven instances of this one defect surfaced while
preparing the paper; the five checks that now fail the build on it, and what each
of them still misses, are in Appendix~\ref{sec:repro}.

\paragraph{Verification moved the numbers in the direction it had to, and changed
no conclusion.} Dropping the four mislabelled pairs raises every phrasing's effect,
as noise centred on zero should: the mean rises from $+0.1917$ to $+0.2075$, the
ratio is unchanged at $4.18\times$ against $4.17\times$, and the effect stays below
$SDC_{95}$. The headline is not an artifact of mislabelled mutants, which was the
most obvious objection to it.

\paragraph{The interaction is real, and we report its size honestly.} Decomposing
the paired effect's variance after main effects cancel: item heterogeneity 0.0410,
arm\,$\times$\,phrasing 0.0046 (sd 0.0678), residual 0.0075. Phrasing is
\textbf{8.7\%} of paired-effect variance -- a modest share of a total dominated by
genuine between-problem heterogeneity, yet large against the effect sizes anyone is
trying to estimate. Quoting only the second half would overstate the case.

\paragraph{Repeatability, the cold path, and call order.} We measured it rather than assuming the usual confound. Repeating a probe 16 times on each of six programs, only the \emph{first} call for a given prompt differs: the remaining 15 are bit-identical, and the seed is inert. So there is no repeatability term in the ordinary sense; \%GRR is reproducibility plus interaction with nothing else hiding in it. What exists instead is a cold-path effect -- a first evaluation differs from every later one by a median of $0.0009$ and at most $0.0115$, which is $0.6\%$ and $7.9\%$ of $SDC_{95}$. Every cell here was measured once and therefore cold, consistently, so the gauge is unbiased by it, but the reported values are cold-prefill values and a pipeline re-querying a cached prompt would see slightly different ones. Single- and repeated-trial designs are not sampling the same quantity. That raises a threat we then tested: the six wordings for a program are called in sequence and share a long prefix, so a cached prefix could make the first cold and the rest warm, and the gauge would book that as appraiser variance. Measuring all six forward and then in reverse -- which moves every wording to a different position -- does shift the readout, by at most $0.0153$. Bounded against what it threatens, the order-induced $SD$ is $0.0015$ against a within-program wording $SD$ of $0.0565$, so it accounts for $0.8\%$ of the appraiser variance. Real, and too small to matter.

\paragraph{Does accuracy rise with stated confidence?} Binning by
confidence, accuracy runs $0.50, 0.63, 0.83, 0.90$ for the first family and
$0.57, 0.70, 0.83, 0.77$ for the second, Spearman $+0.98$ and $+0.83$. The
second's top-bin dip is well inside chance ($p = 0.56$), so the readout is
monotonic in both: ordering by confidence orders by risk within a model.

\paragraph{Wordings or polarities, at a fixed call budget?} The two axes
disagree. Six wordings at one polarity against three at both, matched exactly on
calls and averaged over every choice of subset,
gives \%GRR $37.9$ against $43.9$ and ndc $3.45$ against $2.89$ -- wordings buy
resolution -- while AUROC runs $0.726$ against $0.735$ the other way, a
difference whose interval $[-0.062, +0.047]$ contains zero. Buy wordings, not
polarities, unless ranking is all you need.

\section{Design of the study that was not run}
\label{sec:design}
% ---------------------------------------------------------------------

The apparatus below is built, tested and released; no arm was executed, for
the reason given in Section~\ref{sec:population}.

A seven-arm design over EvalPlus repair
tasks~\citep{liu_evalplus_2023} on HumanEval~\citep{chen_codex_2021} and
MBPP~\citep{austin_mbpp_2021}: a length-matched inert floor, a \emph{task-blind
procedural} placebo generated by the same frozen reflector with every
task-specific field structurally unreachable, two donor-reflection arms, the real
reflection, a compute-matched resample arm, and a known-informative oracle
control. Arms share a byte-identical prompt skeleton; token, item, structural,
call and decoding parity are asserted per item and reported as data rather than
silently repaired; donor assignment is a seeded derangement with a published
manifest; no LLM judge appears anywhere; and the 161-tuple failure pool was frozen
under a single preregistered attempt policy so it could not be redrawn.


\end{document}
